% =====================================================================
%  tikzlib · matematicas/plano-cartesiano
%  Requiere: \usepackage{tikz} ; \usetikzlibrary{arrows.meta}
%  \planocartesiano{xmin}{xmax}{ymin}{ymax}   → plano con rejilla, ejes y números
%  Preview: tikzlib/previews/plano-cartesiano.png
% =====================================================================
\newcommand{\planocartesiano}[4]{%
\begin{tikzpicture}[x=0.7cm, y=0.7cm, font=\sffamily\scriptsize]
  \draw[gray!35, line width=0.4pt] (#1,#3) grid (#2,#4);
  \draw[-{Stealth[length=5pt]}, line width=0.8pt] (#1-0.4,0) -- (#2+0.6,0) node[right]{$x$};
  \draw[-{Stealth[length=5pt]}, line width=0.8pt] (0,#3-0.4) -- (0,#4+0.6) node[above]{$y$};
  \foreach \n in {#1,...,#2}{%
    \ifnum\n=0\else
      \draw[line width=0.6pt] (\n,-0.07) -- (\n,0.07);
      \node[below=1pt, fill=white, inner sep=0.5pt] at (\n,-0.07) {\n};
    \fi}
  \foreach \n in {#3,...,#4}{%
    \ifnum\n=0\else
      \draw[line width=0.6pt] (-0.07,\n) -- (0.07,\n);
      \node[left=1pt, fill=white, inner sep=0.5pt] at (-0.07,\n) {\n};
    \fi}
\end{tikzpicture}}
